% GENERATED FILE. Edit canonical lessons, then run npm run generate:books.
% canonical-source-hash: fnv1a64:5dcd3db8af39b800
% canonical-lessons: GU-C418-sosayti, GU-C418-phlet, GU-C418-banglo, GU-C418-taki, GU-C418-bhadut

\chapter{Societies and Flats}
\label{ch:gu-societies-and-flats}
\hlchaptermodality{418}

\begin{chapteropening}
\textbf{By the end of this chapter:} \emph{I can say where and how I live in a town.}

Complete the last lesson of chapter 418: I can say where and how I live in a town.
\end{chapteropening}

\section[\texorpdfstring{\gu{સોસાયટી} (sosāyṭī)}{સોસાયટી (sosāyṭī)}]{\gu{સોસાયટી} (sosāyṭī) — a housing society; a residential colony}

\label{lesson:GU-C418-sosayti}



\begin{quote}
Before the new one: say the Gujarati for a small gated lane, then the Gujarati for to shut (a door).
\end{quote}

\subsection*{You'll want to know: \gu{સોસાયટી}}
\textbf{\gu{સોસાયટી}} — \emph{sosāyṭī} — \textquotedblleft{}a housing society; a residential colony\textquotedblright{}.

\begin{quote}
\gu{અમારી} \gu{સોસાયટીમાં} \gu{ચાળીસ} \gu{ઘર} \gu{છે}.
\end{quote}

\begin{quote}
\emph{amārī sosāyṭīmā̃ cāḷīs ghar che.}
\end{quote}

\begin{quote}
There are forty houses in our housing society.
\end{quote}

A feminine noun.

\begin{cousinweb}
\textbf{\gu{સોસાયટી}} (\emph{sosāyṭī}) is borrowed from English \emph{society}; in Gujarat a \gu{સોસાયટી} is a planned group of houses with its own gate and committee.
\end{cousinweb}

\begin{grammarlens}[title={What you've built}]
One more word for getting around.
\end{grammarlens}

\subsection*{Your turn}
\begin{itemize}
\raggedright
  \item \emph{Say it:} \emph{sosāyṭī}
  \item \emph{Say it:} \emph{sosāyṭī}, once more
  \item \emph{Say it:} \emph{amārī sosāyṭīmā̃ cāḷīs ghar che.} — There are forty houses in our housing society.
  \item \emph{Recall:} say the Gujarati for a small gated lane, then the Gujarati for to shut (a door), then say \emph{sosāyṭī} again
\end{itemize}

\begin{tcolorbox}[breakable,colback=teal!4,colframe=teal!35!black,title={Before you move on}]
What does \gu{સોસાયટી} mean? (\textquotedblleft{}A housing society; a residential colony\textquotedblright{}.) Say it once more.
\end{tcolorbox}
\section[\texorpdfstring{\gu{ફ્લેટ} (phleṭ)}{ફ્લેટ (phleṭ)}]{\gu{ફ્લેટ} (phleṭ) — a flat (an apartment)}

\label{lesson:GU-C418-phlet}



\begin{quote}
Before the new one: say the Gujarati for to shut (a door), then the Gujarati for a housing society.
\end{quote}

\subsection*{You'll want to know: \gu{ફ્લેટ}}
\textbf{\gu{ફ્લેટ}} — \emph{phleṭ} — \textquotedblleft{}a flat (an apartment)\textquotedblright{}.

\begin{quote}
\gu{તેમનો} \gu{ફ્લેટ} \gu{ત્રીજા} \gu{માળે} \gu{છે}.
\end{quote}

\begin{quote}
\emph{temno phleṭ trījā māḷe che.}
\end{quote}

\begin{quote}
Their flat is on the third floor.
\end{quote}

A masculine noun.

\begin{cousinweb}
\textbf{\gu{ફ્લેટ}} (\emph{phleṭ}) is borrowed from English \emph{flat}, a home on one floor of a building.
\end{cousinweb}

\begin{grammarlens}[title={What you've built}]
One more word for getting around.
\end{grammarlens}

\subsection*{Your turn}
\begin{itemize}
\raggedright
  \item \emph{Say it:} \emph{phleṭ}
  \item \emph{Say it:} \emph{phleṭ}, once more
  \item \emph{Say it:} \emph{temno phleṭ trījā māḷe che.} — Their flat is on the third floor.
  \item \emph{Recall:} say the Gujarati for to shut (a door), then the Gujarati for a housing society, then say \emph{phleṭ} again
\end{itemize}

\begin{tcolorbox}[breakable,colback=teal!4,colframe=teal!35!black,title={Before you move on}]
What does \gu{ફ્લેટ} mean? (\textquotedblleft{}A flat (an apartment)\textquotedblright{}.) Say it once more.
\end{tcolorbox}
\section[\texorpdfstring{\gu{બંગલો} (banglo)}{બંગલો (banglo)}]{\gu{બંગલો} (banglo) — a bungalow; a large house}

\label{lesson:GU-C418-banglo}



\begin{quote}
Before the new one: say the Gujarati for a housing society, then the Gujarati for a flat.
\end{quote}

\subsection*{You'll want to know: \gu{બંગલો}}
\textbf{\gu{બંગલો}} — \emph{banglo} — \textquotedblleft{}a bungalow; a large house\textquotedblright{}.

\begin{quote}
\gu{નદી} \gu{કિનારે} \gu{એક} \gu{સુંદર} \gu{બંગલો} \gu{છે}.
\end{quote}

\begin{quote}
\emph{nadī kināre ek sundar banglo che.}
\end{quote}

\begin{quote}
There is a beautiful bungalow on the river bank.
\end{quote}

A masculine noun.

\begin{cousinweb}
\textbf{\gu{બંગલો}} (\emph{banglo}) is borrowed from English \emph{bungalow}, a word English itself took from India: it is usually traced to Hindi \emph{baṅglā}, \textquotedblleft{}in the Bengal style\textquotedblright{}.
\end{cousinweb}

\begin{grammarlens}[title={What you've built}]
One more word for getting around.
\end{grammarlens}

\subsection*{Your turn}
\begin{itemize}
\raggedright
  \item \emph{Say it:} \emph{banglo}
  \item \emph{Say it:} \emph{banglo}, once more
  \item \emph{Say it:} \emph{nadī kināre ek sundar banglo che.} — There is a beautiful bungalow on the river bank.
  \item \emph{Recall:} say the Gujarati for a housing society, then the Gujarati for a flat, then say \emph{banglo} again
\end{itemize}

\begin{tcolorbox}[breakable,colback=teal!4,colframe=teal!35!black,title={Before you move on}]
What does \gu{બંગલો} mean? (\textquotedblleft{}A bungalow; a large house\textquotedblright{}.) Say it once more.
\end{tcolorbox}
\section[\texorpdfstring{\gu{ટાંકી} (ṭā̃kī)}{ટાંકી (ṭā̃kī)}]{\gu{ટાંકી} (ṭā̃kī) — a water tank}

\label{lesson:GU-C418-taki}



\begin{quote}
Before the new one: say the Gujarati for a flat, then the Gujarati for a bungalow.
\end{quote}

\subsection*{You'll want to know: \gu{ટાંકી}}
\textbf{\gu{ટાંકી}} — \emph{ṭā̃kī} — \textquotedblleft{}a water tank\textquotedblright{}.

\begin{quote}
\gu{છત} \gu{પરની} \gu{ટાંકીમાં} \gu{પાણી} \gu{નથી}.
\end{quote}

\begin{quote}
\emph{chat parnī ṭā̃kīmā̃ pāṇī nathī.}
\end{quote}

\begin{quote}
There is no water in the tank on the roof.
\end{quote}

A feminine noun.

\begin{cousinweb}
\textbf{\gu{ટાંકી}} (\emph{ṭā̃kī}) goes back to \emph{*ṭaṅka}, a form CDIAL reconstructs; CDIAL lists the Gujarati word, \textquotedblleft{}a cistern\textquotedblright{} (CDIAL 5430).
\end{cousinweb}

\begin{grammarlens}[title={What you've built}]
One more word for getting around.
\end{grammarlens}

\subsection*{Your turn}
\begin{itemize}
\raggedright
  \item \emph{Say it:} \emph{ṭā̃kī}
  \item \emph{Say it:} \emph{ṭā̃kī}, once more
  \item \emph{Say it:} \emph{chat parnī ṭā̃kīmā̃ pāṇī nathī.} — There is no water in the tank on the roof.
  \item \emph{Recall:} say the Gujarati for a flat, then the Gujarati for a bungalow, then say \emph{ṭā̃kī} again
\end{itemize}

\begin{tcolorbox}[breakable,colback=teal!4,colframe=teal!35!black,title={Before you move on}]
What does \gu{ટાંકી} mean? (\textquotedblleft{}A water tank\textquotedblright{}.) Say it once more.
\end{tcolorbox}
\section[\texorpdfstring{\gu{ભાડૂત} (bhāḍūt)}{ભાડૂત (bhāḍūt)}]{\gu{ભાડૂત} (bhāḍūt) — a tenant}

\label{lesson:GU-C418-bhadut}



\begin{quote}
Before the new one: say the Gujarati for a bungalow, then the Gujarati for a water tank.
\end{quote}

\subsection*{You'll want to know: \gu{ભાડૂત}}
\textbf{\gu{ભાડૂત}} — \emph{bhāḍūt} — \textquotedblleft{}a tenant\textquotedblright{}.

\begin{quote}
\gu{ભાડૂત} \gu{મહિનાના} \gu{અંતે} \gu{ભાડું} \gu{ચૂકવે} \gu{છે}.
\end{quote}

\begin{quote}
\emph{bhāḍūt mahinānā ante bhāḍũ cūkve che.}
\end{quote}

\begin{quote}
The tenant pays the rent at the end of the month.
\end{quote}

A masculine or feminine noun, by the person meant.

\begin{cousinweb}
\textbf{\gu{ભાડૂત}} (\emph{bhāḍūt}) is built on \gu{ભાડું}, rent, which CDIAL lists under \emph{*bhārta}, \textquotedblleft{}pertaining to servants\textquotedblright{} (CDIAL 9468): the one who pays the rent.
\end{cousinweb}

\begin{grammarlens}[title={What you've built}]
One more word for getting around.
\end{grammarlens}

\subsection*{Your turn}
\begin{itemize}
\raggedright
  \item \emph{Say it:} \emph{bhāḍūt}
  \item \emph{Say it:} \emph{bhāḍūt}, once more
  \item \emph{Say it:} \emph{bhāḍūt mahinānā ante bhāḍũ cūkve che.} — The tenant pays the rent at the end of the month.
  \item \emph{Recall:} say the Gujarati for a bungalow, then the Gujarati for a water tank, then say \emph{bhāḍūt} again
\end{itemize}

\begin{tcolorbox}[breakable,colback=teal!4,colframe=teal!35!black,title={Before you move on}]
What does \gu{ભાડૂત} mean? (\textquotedblleft{}A tenant\textquotedblright{}.) Say it once more.
\end{tcolorbox}
